\section{Conclusion}\label{sec:conclusion}
In this paper, we propose the \method, a simple yet effective plug-in design for various kinds of long-tail semantic segmentation tasks. 
%
We introduce category scale-related variation during the model training stage.
%
This variation is inversely proportional to the frequency of occurrences of instances, which effectively closes the gap between the feature area of different categories.
%
Extensive experiments on fully supervised, semi-supervised, and unsupervised domain adaptive semantic segmentation tasks suggest our method can boost performance.
%
Compared with other methods for alleviating the class-imbalance issues, our \method is better and more concise and general.
%
Furthermore, sufficient ablation experiments as well as intuitive visualization results prove the necessity of individual components and the effectiveness of our method.
%
%\method has the potential to be extended to other tasks and become a general paradigm.


\noindent\textbf{Discussion.} One necessary premise of \method is that the category frequencies need to be known.
%
It is unlikely to be satisfied in some tasks like unsupervised semantic segmentation and  domain generalized semantic segmentation.
%
% Besides, this method may involve privacy issues of training data.

% Although this premise can be approximated by the source domain data category distribution in UDA tasks,
%